\subsubsection*{Solution}
The four observations are mutually compatible, but only if the magnetic order carries a uniform component in addition to the two reported propagation vectors. The two neutron peaks and the birefringence alone fix a magnetic cell and a point group; the Kerr effect then rules out every group that those two constraints would otherwise permit, and forces the remaining symmetry down to a single maximal magnetic space group.

\paragraph*{Step-by-Step Derivation}

\subparagraph*{1. Parent symmetry and the magnetic site}

Space group \#182 is $P6_322$, of point group $622$, with twelve operations per primitive cell. All twelve are proper rotations: a $6_3$ screw along $c$ and its powers, and six two-fold axes in the plane, three at $z=0$ along $\langle 100\rangle$ and three at $z=1/4$ along $\langle 1\bar{1}0\rangle$. There is no inversion centre, no mirror and no rotoinversion, so the crystal is chiral. The whole group is generated by two operations together with the lattice translations,

\[
\{6_3^{+}\,|\,0,0,\tfrac12\}:(x-y,\,x,\,z+\tfrac12), \qquad \{2_{[110]}\,|\,0\}:(y,\,x,\,-z).
\]

Wyckoff position $2c$ is $(\tfrac13,\tfrac23,\tfrac14)$ and $(\tfrac23,\tfrac13,\tfrac34)$, with site symmetry $32$: the site lies where the three-fold axis along $c$ crosses the three in-plane two-folds at $z=1/4$. The two sites form triangular layers at $z=c/4$ and $z=3c/4$, exchanged by the $6_3$ screw. This is the structure type of the intercalated dichalcogenides $M_{1/3}\mathrm{NbS}_2$, of which edgarite $\mathrm{FeNb_3S_6}$ is the mineral (\href{\detokenize{https://doi.org/10.1016/0022-3697(70)90315-X}}{J.\ Phys.\ Chem.\ Solids \textbf{31}, 1057 (1970)}).

\subparagraph*{2. The two peaks fix the magnetic cell}

Magnetic Bragg peaks occur at $\mathbf{Q}=\mathbf{G}\pm\mathbf{k}$ with $\mathbf{G}$ a parent reciprocal lattice vector, so the measured peak positions are the propagation vectors. Acting with $622$ on $(\tfrac12,0,0)$ generates the three-arm star of $M$,

\[
\{\,(\tfrac12,0,0),\ (0,\tfrac12,0),\ (\tfrac12,\tfrac12,0)\equiv(-\tfrac12,\tfrac12,0)\,\},
\]

so $\mathbf{q}_1$ and $\mathbf{q}_2$ are two of the three arms and the third, $\mathbf{q}_3=(\tfrac12,0,0)$, carries no observed intensity. A parent translation $\mathbf{t}$ leaves the structure invariant iff $e^{-i\mathbf{k}\cdot\mathbf{t}}=1$ for every $\mathbf{k}$ present. With $\mathbf{t}=n_1\mathbf{a}_1+n_2\mathbf{a}_2+n_3\mathbf{a}_3$ the two conditions read $n_2/2\in\mathbb{Z}$ and $(n_1+n_2)/2\in\mathbb{Z}$, forcing $n_1$ and $n_2$ both even. Hence

\[
{\text{magnetic lattice}=\{2\mathbf{a}_1,\,2\mathbf{a}_2,\,\mathbf{a}_3\},\qquad \text{a } 2a\times 2b\times c \text{ cell of index } 4,\ 8 \text{ magnetic atoms.}}
\]

Expand the out-of-plane moment on layer $s$ in the cell $(n_1,n_2)$ as a Fourier sum over $\Gamma$ and the three arms. Because $2\mathbf{k}$ is a reciprocal lattice vector for each arm, every phase $e^{-i\mathbf{k}\cdot\mathbf{R}}$ is $\pm1$ and every amplitude is real:

\[
m_z(s,n_1,n_2)=c_{s,0}+c_{s,1}(-1)^{n_2}+c_{s,2}(-1)^{n_1+n_2}+c_{s,3}(-1)^{n_1}.
\]

The reported peaks make $c_{s,1}$ and $c_{s,2}$ non-zero and set $c_{s,3}=0$. They say nothing about the uniform term $c_{s,0}$: a $k=0$ component overlaps the nuclear Bragg reflections instead of producing a superlattice one, so it is not constrained by the quoted peak positions and may simply be unreported. It is therefore carried along:

\[
m_z(s,n_1,n_2)=c_{s,0}+c_{s,1}(-1)^{n_2}+c_{s,2}(-1)^{n_1+n_2}.
\]

\subparagraph*{3. Birefringence removes the trigonal and hexagonal candidates}

A hexagonal crystal is uniaxial and therefore already birefringent for a general propagation direction; what the transition can add is birefringence for light along $c$, where the parent shows none. A three-fold or six-fold axis about $c$ forces the in-plane block of the dielectric tensor to be isotropic, $\epsilon_{xx}=\epsilon_{yy}$ and $\epsilon_{xy}=0$, so in-plane birefringence requires that axis to be absent. Any structure using only two arms breaks it automatically, since a three-fold rotation permutes the arms cyclically. The condition therefore excludes only the equal-amplitude three-arm candidates, BNS $149.23$, $182.182$ and $182.183$, which keep the axis. Three-state nematicity of exactly this kind is observed optically in $\mathrm{Fe_{1/3}NbS_2}$ (\href{\detokenize{https://www.nature.com/articles/s41563-020-0681-0}}{Nature Materials \textbf{19}, 1062 (2020)}).

\subparagraph*{4. The Kerr effect is the decisive constraint}

A Kerr rotation requires a non-zero antisymmetric part of the dielectric tensor, a time-odd axial vector, which is allowed only if the magnetic point group is one of the $31$ permitting a spontaneous magnetization. This excludes every type-IV magnetic space group, since an anti-translation $\{1'|\mathbf{t}\}$ sends any macroscopic time-odd quantity to minus itself while leaving the crystal alone, forcing it to vanish.

Apply the translation $\mathbf{a}_2$, that is $n_2\to n_2+1$, to the moment of step 2. Both modulated terms change sign while the uniform term does not:

\[
m_z(s,n_1,n_2+1)=c_{s,0}-\left[c_{s,1}(-1)^{n_2}+c_{s,2}(-1)^{n_1+n_2}\right]=2c_{s,0}-m_z(s,n_1,n_2).
\]

So $\{1'|\mathbf{a}_2\}$ reverses every moment, and the group is type IV, precisely when $c_{s,0}=0$. The two-arm modulation on its own cancels under that translation for any amplitudes, and an explicit enumeration of every such pattern returns only the type-IV groups $C_a222_1$ $(20.36)$, $C_a2$ $(5.17)$, $P_a2_1$ $(4.10)$ and $P_S1$ $(1.3)$, none of which permits a Kerr rotation. The observed Kerr effect therefore requires

\[
{ c_{s,0}\neq 0: \text{ a uniform component along } c \text{ must be present.} }
\]

Equivalently, the three arms carry the three non-trivial characters of $\mathbb{Z}_2\times\mathbb{Z}_2$, and for any two of them exactly one translation makes both characters $-1$; only the uniform term is blind to it. In Landau language the lowest invariant coupling the uniform magnetization to the $M$-point order parameters is $M_z\,\eta_1\eta_2\eta_3$, which is quartic and needs all three arms, so no invariant $M_z\eta_i\eta_j$ exists.

\subparagraph*{5. The nature of the uniform component}

Because the order is out of plane, $c_{s,0}\neq0$ is a ferrimagnetic imbalance rather than a canting: the up and down sites simply do not carry equal moments. An equal-moment Ising double-$q$ state does not exist in any case, since $|c_1+c_2|=|c_1-c_2|$ forces $c_1=0$ or $c_2=0$, which is single-$q$. In an unpolarized diffraction measurement the effect of such a component on the nuclear peak intensities is quadratic in the moment and easily below systematic error, which is consistent with its not appearing in the reported peak list.

\subparagraph*{6. The surviving symmetry}

Only operations mapping the pair $\{\mathbf{q}_1,\mathbf{q}_2\}$ onto itself can survive, that is the little co-group of the unused arm: the identity, the $2_1$ screw along $c$, and two mutually perpendicular in-plane two-folds. With the moments along $c$ the in-plane two-folds reverse $m_z$ and must be primed, while the $c$-axis screw stays unprimed. The magnetic point group is $2'2'2$ with the unprimed axis along $c$, which permits $M\parallel c$, and the magnetic space group is $C2'2'2_1$ in the $C$-centred orthorhombic magnetic cell $a_m=2a$, $b_m=2\sqrt{3}\,a$, $c_m=c$. Its index-2 subgroups that still permit $M\parallel c$ are $C2'$, $P2_1$ and $P1$; the remaining maximal candidate $C22'2_1'$ ($20.34$) is excluded because its symmetry-allowed moment lies in the plane and so produces no polar Kerr signal for out-of-plane order.

\[
{ \text{BNS } 20.33=C2'2'2_1\ \text{(maximal)},\qquad 5.15=C2',\qquad 4.7=P2_1,\qquad 1.1=P1. }
\]

All group identifications were made with \texttt{spglib}'s magnetic space group database and cross-checked against \texttt{pymatgen}.

\subparagraph*{Remark}

The vanishing of the third arm is not enforced by $20.33$, which permits a $(\tfrac12,0,0)$ component. No magnetic space group can simultaneously forbid the third peak and permit a Kerr rotation, because forbidding that arm requires the very anti-translation that kills the Kerr effect. If the third peak is instead merely weak and unreported, the structure is a three-arm state with unequal amplitudes and the answer is unchanged.
